По замечанию \ref{remUltraMetric} в облаке \( [\Delta _1] \) для
любых пространств \( X_1, X_2 \) выполняется ультраметрическое
неравенство. Целью данного раздела является показать, что в общем
случае это не верно.
\begin{lemma}
  Если $X$ --- подмножество вещественной прямой. Положим
  \[
    d =\sup\{r\in \R_+ \colon \exists a, (a-r, a+r)\subset \R
    \setminus X \}.
  \]
  Тогда \( d_{GH}(X,\R) \) равно \( d \).
  \label{lemmaDiamDist}
\end{lemma}
\begin{proof} Каноническое вложение \( X \hookrightarrow \R \) задает
  реализацию с расстоянием Хаусдорфа равным \( d \).
  Отсюда
  \begin{equation}
    d_{GH}(X, \R) \le d.
    \label{eq:lemmaDiamDist}
  \end{equation}

  Далее считаем, что в $\mathbb{R} \setminus X$ лежит
  интервал $(a-r, a+r)$, \( r \le d \).
  Предположим, что $d_{G H}(\mathbb{R}, X)<r$. Тогда должна
  существовать реализация
  $\left(\mathbb{R}^{\prime}, X^{\prime}, Y\right)$ такая, что
  $d_H\left(\mathbb{R}^{\prime},
  X^{\prime}\right)=$ $r^{\prime}<r$. Обозначим
  \[
    X_l \coloneqq \left\{x\in X \colon x \le a-r\right\},
    X_r \coloneqq \left\{x\in X \colon x \ge a+r\right\}.
  \]
  Это подмножества \( X \) лежащие слева и справа от интервала \(
  (a-r, a+r) \) соответственно. Заметим, что \( X = X_l \sqcup X_r
  \). Также обозначим \( X_l', X_r'\) их образы в реализации.
  Далее обозначим \[U_l:=\cup_{x \in
    X_l'} B\left(x,
    r^{\prime}+\frac{r-r^{\prime}}{2}\right),~~ U_r:=\cup_{x \in
    X_r'} B\left(x,
  r^{\prime}+\frac{r-r^{\prime}}{2}\right),\]
  то есть $U_l \cup U_r$ --- покрытие $X^{\prime}$ шарами радиуса
  $r^\prime + \frac{r-r^{\prime}}{2}$.
  Получаем, что $U_l, U_r$ - два открытых непересекающихся множества,
  но также $\mathbb{R}^{\prime} \in U_l \cup U_r$ по определению
  расстояния Хаусдорфа, что противоречит
  связности прямой.

  Получили, что \( d_{GH}(X. \R) \ge r \) для всех \( r \) таких, что
  в \( \R \setminus X \) лежит интервал диаметра \( 2r \). Переходя к
  \( \sup \) получаем \( d_{GH}(X, \R) \ge d \). Вместе с
  неравенством \ref{eq:lemmaDiamDist} получаем искомое равенство \(
  d_{GH}(X, \R) = d\).
\end{proof}
Нам понадобится следующая вспомогательная лемма, аналогичная лемме
\ref{lemDiamImage}.
\begin{lemma}
  Если \( R \) - соответствие между метрическими пространствами \( X
  \) и \( Y \), то \( \diam R(x) \le \dis R \) для любой точки \( x\in X \).
\end{lemma}
Обозначим \( \widetilde \R \) подпространство \( \R^2 \) с метрикой
\( l_1 \), состоящее из прямой \( \big\{(x,0) \colon x\in \R\big\} \)
и точки \( (0,1) \).
\begin{theorem}
  Пусть \( \Z \) --- метрическое  пространство целых чисел с метрикой
  из \( \R \). Тогда
  \begin{enumerate}
    \item \( d_{GH}(\Z, \R) = \frac 1 2 \), \( d_{GH}(\widetilde \R,
      \R)\le \frac 1 2 , \)\label{thrmPt:1}
    \item \( d_{GH}(\Z, \widetilde \R) \ge \frac 2 3 . \)\label{thrmPt:2}
  \end{enumerate}
  \label{thrmRUltraMetric}
\end{theorem}
\begin{proof}
  (\ref{thrmPt:1}). Каноничное вложение \( \Z \hookrightarrow \R \)
  задает реализацию с расстоянием Хаусдорфа равным \( \frac 1 2 \).
  Отсюда \( d_{GH}(\Z,\R) \le \frac 1 2 \). Обратное неравенство
  следует из леммы \ref{lemmaDiamDist}. Далее, вложение \( \widetilde
  \R \subset \R^2 \) и \( \R' = \left\{(x,\frac 1 2) \colon x \in \R\right\} \)
  дает реализацию с расстоянием Хаусдорфа равным \( \frac 1 2 \).

  (\ref{thrmPt:2}).
  Пусть $R$ --- соответствие
  между $\mathbb{Z}$ и $\widetilde{\mathbb{R}}$, с искажением,
  равным $1 + \varepsilon$ и
  в образе точки $i$ из $\mathbb{Z}$ лежит $(0,1)$. По лемме
  \ref{lemDiamImage} диаметр
  образа точки не может быть больше искажения соответствия,
  следовательно, образ $i$ лежит в $(-\varepsilon, \varepsilon) \cup
  \bigl\{(0,1)\bigr\}$. Это
  означает, что для
  $x$ не лежащих в $(-\varepsilon, \varepsilon)$, пара $(i, x)$ не лежит в $R$.
  Обозначим через $\mathcal{N}$ множество всех целых чисел таких, что их
  образ лежит в $(-\varepsilon, \varepsilon) \cup
  \bigl\{(0,1)\bigr\}$. Множество
  $\mathcal{N}$ не пусто и не равно $\mathbb{Z}$, следовательно, по
  лемме \ref{lemmaDiamDist},
  расстояние от $\mathbb{Z} \setminus\mathcal{N}$ до $\mathbb{R}$ будет не
  меньше 1.
  Из соответствия $R$ уберем пару $\bigl(i,(0,1)\bigr)$, а также
  все пары $(k,x)$ такие, что $x \in  (-\varepsilon, \varepsilon)$.
  Получившееся множество обозначим $R'$.
  Так как все точки из $\mathbb{R}\setminus(-\varepsilon,
  \varepsilon)$ лежат в $R$
  только в паре с точками из $ \mathbb{Z} \setminus\mathcal{N}$ и наоборот,
  множество $R'$ будет соответствием между
  $\mathbb{R}\setminus(-\varepsilon, \varepsilon)$ и
  $ \mathbb{Z} \setminus\mathcal{N}$. Искажение
  подмножества соответствия по определению не больше искажения
  самого соответствия. Получаем цепочку неравенств:
  \[1+\varepsilon = \dis R \ge \dis R' \ge
    2d_{GH}\bigl(\mathbb{R}\setminus(-\varepsilon, \varepsilon), \mathbb{Z}
  \setminus \mathcal{N}\bigr).\]
  По неравенству треугольника
  \[
    2d_{GH}\bigl(\mathbb{R}\setminus(-\varepsilon, \varepsilon), \mathbb{Z}
    \setminus \mathcal{N}\bigr) \ge
    2\Big|d_{GH}\bigl(\mathbb{R}, \mathbb{Z} \setminus
    \mathcal{N}\bigr) - d_{GH}\bigl(\mathbb{R}\setminus(-\varepsilon,
    \varepsilon), \mathbb{R}\bigr)\Big| \ge 2 - 2\varepsilon.
  \]
  Получили неравенство: $1 + \varepsilon \ge 2 - 2\varepsilon$.
  Из него получаем нижнюю оценку на $\varepsilon$:
  \[\varepsilon \ge \frac 1 3,\]
  откуда $\dis R \ge \frac 4 3$ и
  $d_{GH}\left(\widetilde{\mathbb{R}}, \mathbb{Z}\right) \ge \frac 2
  3 $.
\end{proof}
